\section{Rerooting DP}
\whenuse{Se necesita una respuesta para cada vértice considerando que ese
vértice es la raíz o punto de inicio: suma de distancias a todos, cantidad de
nodos alcanzables, mejor costo al elegir una capital o valor de una DP para
cada raíz. La respuesta de raíces vecinas debe poder actualizarse reutilizando
la información de sus subárboles. Evita ejecutar una DFS $O(n)$ desde cada
vértice, lo que costaría $O(n^2)$.}

\complexity{tiempo $O(n)$; memoria $O(n)$.}

\begin{lstlisting}[style=cpp]
struct DistanceRerooting {
    int node_count;
    const vector<vector<int>> &tree;
    vector<int> subtree_size;
    vector<ll> distance_sum;

    explicit DistanceRerooting(const vector<vector<int>> &input_tree)
        : node_count(input_tree.size()), tree(input_tree),
          subtree_size(node_count), distance_sum(node_count) {
        dfs_down(0, -1);
        dfs_up(0, -1);
    }

    void dfs_down(int vertex, int parent) {
        subtree_size[vertex] = 1;
        for (int to : tree[vertex]) if (to != parent) {
            dfs_down(to, vertex);
            subtree_size[vertex] += subtree_size[to];
            distance_sum[vertex] += distance_sum[to] + subtree_size[to];
        }
    }
    void dfs_up(int vertex, int parent) {
        for (int to : tree[vertex]) if (to != parent) {
            distance_sum[to] = distance_sum[vertex]
                             - subtree_size[to] + (node_count - subtree_size[to]);
            dfs_up(to, vertex);
        }
    }
};
\end{lstlisting}
\begin{lstlisting}[style=cpp]
tree[u].push_back(v);
tree[v].push_back(u);

DistanceRerooting rerooting(tree);
for(int v=0;v<n;++v)
    cout<<rerooting.distance_sum[v]<<'\n';
\end{lstlisting}
Para una combinaci\'on gen\'erica no invertible, usar prefijos y sufijos de las
contribuciones de los hijos para excluir uno en $O(1)$.
